\chapter{Lập trình hướng đối tượng trong Python}

Lập trình hướng đối tượng (Object-Oriented Programming - OOP) là phương pháp luận lập trình đóng vai trò then chốt trong việc thiết kế các thư viện Machine Learning hiện đại như Scikit-Learn, PyTorch và TensorFlow. Chương này trình bày các khái niệm OOP cốt lõi và hướng dẫn cách ứng dụng OOP để tự đóng gói một mô hình Machine Learning.

\section{Lớp và Đối tượng (Class and Object)}

\begin{definition}[Lớp và Đối tượng]
\textbf{Lớp (Class)} là khuôn mẫu định nghĩa các thuộc tính (dữ liệu) và phương thức (hành vi). \textbf{Đối tượng (Object)} là một thể hiện cụ thể được khởi tạo từ lớp đó.
\end{definition}

\begin{example}
Khai báo một lớp cơ bản lưu trữ mẫu dữ liệu:
\end{example}

\begin{lstlisting}[caption={Khai báo Class và tạo Object trong Python}]
class DataSample:
    def __init__(self, feature_name, value):
        self.feature_name = feature_name  # Thuoc tinh
        self.value = value

    def display_info(self):              # Phuong thuc
        print(f"Feature: {self.feature_name}, Value: {self.value}")

# Khoi tao doi tuong
sample = DataSample(feature_name="Age", value=25)
sample.display_info()
\end{lstlisting}

\section{Các phương thức đặc biệt (Magic Methods)}

Các phương thức có cú pháp hai dấu gạch dưới (\texttt{\_\_method\_\_}) giúp tùy biến hành vi của đối tượng trong Python.

\subsection{Phương thức \texttt{\_\_init\_\_} và \texttt{\_\_call\_\_}}
\begin{itemize}
    \item \texttt{\_\_init\_\_}: Hàm khởi tạo (constructor), tự động thực thi khi tạo mới một đối tượng.
    \item \texttt{\_\_call\_\_}: Cho phép gọi trực tiếp một đối tượng giống như một hàm, thường dùng khi thiết kế các Layer trong Neural Network.
\end{itemize}

\begin{lstlisting}[caption={Sử dụng phương thức \_\_call\_\_}]
class Multiplier:
    def __init__(self, factor):
        self.factor = factor

    def __call__(self, x):
        return x * self.factor

# Su dung nhu mot ham
double = Multiplier(factor=2)
print(double(5))  # Output: 10
\end{lstlisting}

\section{Tính Kế thừa (Inheritance)}

Tính kế thừa cho phép tạo một lớp mới dựa trên lớp đã có, giúp tái sử dụng mã nguồn và chuẩn hóa giao diện các mô hình.

\begin{lstlisting}[caption={Kế thừa lớp mô hình cơ sở}]
class BaseModel:
    def __init__(self, name):
        self.name = name

    def fit(self, X, y):
        raise NotImplementedError("Phuong thuc fit() phai duoc cai dat o lop con.")

class LinearModel(BaseModel):
    def __init__(self):
        super().__init__(name="Linear Regression")

    def fit(self, X, y):
        print(f"Dang huan luyen mo hinh {self.name}...")

# Khoi tao va huan luyen
model = LinearModel()
model.fit(X=None, y=None)
\end{lstlisting}

\section{Thiết kế Custom Estimator chuẩn Scikit-Learn}

Trọng tâm ứng dụng OOP trong Machine Learning là đóng gói thuật toán thành một Class có các giao diện chuẩn (\texttt{fit}, \texttt{predict}).

\begin{lstlisting}[caption={Đóng gói thuật toán thành Class trong Python}]
import numpy as np

class CustomLinearRegressor:
    def __init__(self, learning_rate=0.01, epochs=1000):
        self.lr = learning_rate
        self.epochs = epochs
        self.weights = None
        self.bias = None

    def fit(self, X, y):
        n_samples, n_features = X.shape
        self.weights = np.zeros(n_features)
        self.bias = 0.0

        for _ in range(self.epochs):
            y_predicted = np.dot(X, self.weights) + self.bias
            dw = (1 / n_samples) * np.dot(X.T, (y_predicted - y))
            db = (1 / n_samples) * np.sum(y_predicted - y)

            self.weights -= self.lr * dw
            self.bias -= self.lr * db

    def predict(self, X):
        return np.dot(X, self.weights) + self.bias
\end{lstlisting}

\section{Tóm tắt chương}

\begin{itemize}
    \item Khái niệm Lớp, Đối tượng và tư duy lập trình hướng đối tượng.
    \item Ứng dụng các Magic Methods (\texttt{\_\_init\_\_}, \texttt{\_\_call\_\_}) trong xây dựng Framework.
    \item Nguyên lý Kế thừa giúp tái sử dụng và quản lý mã nguồn linh hoạt.
    \item Đóng gói thuật toán Machine Learning thành Class theo giao diện tiêu chuẩn.
\end{itemize}